%
% 000-session19.tex
% Presentation Session 19: Practice 7 -- JSON REST API (Ensembl)
%
% Compile to .pdf with LaTeX (pdflatex)
% Requires Beamer package (latex-beamer)
%

\documentclass[xcolor=table]{beamer}
\usetheme{Warsaw}
\beamertemplatenavigationsymbolsempty
\setbeamertemplate{headline}{}
\setbeamerfont{title}{size=\large}
\useoutertheme{infolines}
\usepackage[english]{babel}
\usepackage[utf8]{inputenc}
\usepackage{graphics}
\usepackage{amssymb}
\usepackage{multirow}
\usepackage{xcolor}
\usepackage{framed}
\usepackage{hyperref}
\usepackage{tikz}
\usetikzlibrary{arrows.meta, positioning}

\definecolor{shadecolor}{RGB}{180,180,180}

%%%%%%%%%%%%%%%%%%%%%%%%%%%%%%%%%%%%%%%%%%%%%%%%%%%%%%%%%%%%%%%%
% Python Highlighting
%%%%%%%%%%%%%%%%%%%%%%%%%%%%%%%%%%%%%%%%%%%%%%%%%%%%%%%%%%%%%%%%
\DeclareFixedFont{\ttb}{T1}{txtt}{bx}{n}{10}
\DeclareFixedFont{\ttm}{T1}{txtt}{m}{n}{10}

% Custom colors
\usepackage{color}
\definecolor{deepblue}{rgb}{0,0,0.5}
\definecolor{deepred}{rgb}{0.6,0,0}
\definecolor{deepgreen}{rgb}{0,0.5,0}
\definecolor{grey1}{rgb}{0.5,0.5,0.5}

\usepackage{listings}

% Python style for highlighting
\newcommand\pythonstyle{\lstset{
language=Python,
basicstyle=\ttm\small,
morekeywords={self},
keywordstyle=\ttb\color{deepblue},
emphstyle=\ttb\color{deepred},
stringstyle=\color{deepgreen},
commentstyle=\small\color{grey1}\ttm,
frame=single,
showstringspaces=false
}}

% HTML style for highlighting
\newcommand\htmlstyle{\lstset{
language=HTML,
basicstyle=\ttm\small,
keywordstyle=\ttb\color{deepblue},
stringstyle=\color{deepgreen},
commentstyle=\small\color{grey1}\ttm,
frame=single,
showstringspaces=false
}}

% Python environment
\lstnewenvironment{python}[1][]
{
\pythonstyle
\lstset{#1}
}
{}

% HTML environment
\lstnewenvironment{htmlcode}[1][]
{
\htmlstyle
\lstset{#1}
}
{}

% Python for inline
\newcommand\pythoninline[1]{{\pythonstyle\lstinline!#1!}}

% HTML for inline
\newcommand\htmlinline[1]{{\htmlstyle\lstinline!#1!}}

\newcommand{\asignatura}{Session 19: Practice 7 -- JSON REST API}
\newcommand{\grado}{Programming in Network Environments}
\newcommand{\curso}{2025-2026}

\title[\asignatura]{\asignatura}
\subtitle{\grado}
\author{URJC}
\institute{GSyC}
\date{Course \curso}

\AtBeginSection[]
{
\begin{frame}<beamer>
\begin{center}
{\LARGE \insertsection}
\end{center}
\end{frame}
}

\begin{document}

\frame{
\maketitle
}

\begin{frame}
\frametitle{Contents}
\tableofcontents
\end{frame}

%%%%%%%%%%%%%%%%%%%%%%%%%%%%%%%%%%%%%%%%%%%%%%%%%%%%%%%%%%%%%%%%
%%%%%%%%%%%%%%%%%%%%%%%%%%%%%%%%%%%%%%%%%%%%%%%%%%%%%%%%%%%%%%%%
% SECTION 1: Session 17 Review
%%%%%%%%%%%%%%%%%%%%%%%%%%%%%%%%%%%%%%%%%%%%%%%%%%%%%%%%%%%%%%%%
%%%%%%%%%%%%%%%%%%%%%%%%%%%%%%%%%%%%%%%%%%%%%%%%%%%%%%%%%%%%%%%%

\section{Review: Session 17 -- Seq2 Server}

\begin{frame}[fragile]
\frametitle{What We Learned in Session 17 (1/2)}

\textbf{Seq2 Server}: a web server using \pythoninline{http.server} with 4 services.

\begin{itemize}
\item \textbf{PING}: test if the server is alive (\texttt{/ping})
\vspace{0.15cm}
\item \textbf{GET}: retrieve a sequence by number (\texttt{/get?n=2})
\vspace{0.15cm}
\item \textbf{GENE}: get a gene's complete sequence (\texttt{/gene?name=U5})
\vspace{0.15cm}
\item \textbf{OPERATION}: perform operations on a sequence (\texttt{info}, \texttt{comp}, \texttt{rev})
\end{itemize}

\end{frame}

\begin{frame}[fragile]
\frametitle{What We Learned in Session 17 (2/2)}

\textbf{Key concepts}:

\begin{itemize}
\item HTML \textbf{forms} to send data to the server
\item Server-side \textbf{URL parsing} with \pythoninline{urlparse} and \pythoninline{parse_qs}
\item \textbf{Dynamic HTML} generation with \pythoninline{jinja2} templates
\item \textbf{Error handling}: return error pages for unknown resources
\end{itemize}

\vspace{0.3cm}

\textbf{Server structure}:

\begin{python}
class Seq2Handler(
        http.server.BaseHTTPRequestHandler):
    def do_GET(self):
        url_path = urlparse(self.path)
        path = url_path.path
        arguments = parse_qs(url_path.query)
\end{python}

\end{frame}

%%%%%%%%%%%%%%%%%%%%%%%%%%%%%%%%%%%%%%%%%%%%%%%%%%%%%%%%%%%%%%%%
%%%%%%%%%%%%%%%%%%%%%%%%%%%%%%%%%%%%%%%%%%%%%%%%%%%%%%%%%%%%%%%%
% SECTION 2: Session 19 -- Practice 7
%%%%%%%%%%%%%%%%%%%%%%%%%%%%%%%%%%%%%%%%%%%%%%%%%%%%%%%%%%%%%%%%
%%%%%%%%%%%%%%%%%%%%%%%%%%%%%%%%%%%%%%%%%%%%%%%%%%%%%%%%%%%%%%%%

\section{Session 19: Practice 7 -- JSON REST API}

\begin{frame}
\frametitle{Session 19 Goals}

\textbf{In this session we will}:

\begin{itemize}
\item Practice with \textbf{JSON} and the \textbf{API REST}
\item Read information from the \textbf{Ensembl project} database
\item Develop a \textbf{python client} that connects to a remote REST API
\item Perform \textbf{calculations} on gene sequences retrieved from Ensembl
\end{itemize}

\vspace{0.4cm}

\textbf{Time}: 2 hours

\vspace{0.2cm}

\textbf{Folder}: \texttt{P07/}

\vspace{0.2cm}

\textbf{Important}: Finish all S18 exercises before starting!

\end{frame}

\begin{frame}
\frametitle{The Ensembl Database}

The \href{http://www.ensembl.org/index.html}{\textbf{Ensembl}} database can be accessed in \textbf{two ways}:

\vspace{0.3cm}

\begin{enumerate}
\item \textbf{Browser}: write the URL in the address bar $\rightarrow$ web page for \textbf{humans}
\vspace{0.2cm}
\item \textbf{API REST}: applications get objects directly from the server $\rightarrow$ interface for \textbf{programs}
\end{enumerate}

\vspace{0.3cm}

The API REST documentation:
\begin{center}
\url{https://rest.ensembl.org/}
\end{center}

\vspace{0.2cm}

Resources are accessed through \textbf{endpoints} (named paths).

\end{frame}

\begin{frame}
\frametitle{Genes to Analyze}

The application should read \textbf{10 genes} and perform calculations on them:

\vspace{0.3cm}

\begingroup\small
\begin{center}
\begin{tabular}{|l|l|}
\hline
\textbf{Gene} & \textbf{Gene} \\
\hline
FRAT1 & TTTY4C \\
\hline
ADA & RBMY2YP \\
\hline
FXN & FGFR3 \\
\hline
RNU6\_269P & KDR \\
\hline
MIR633 & ANK2 \\
\hline
\end{tabular}
\end{center}
\endgroup

\vspace{0.3cm}

\textbf{Calculations}: total bases, count of each base, percentage, most frequent base.

\end{frame}

%%%%%%%%%%%%%%%%%%%%%%%%%%%%%%%%%%%%%%%%%%%%%%%%%%%%%%%%%%%%%%%%
%%%%%%%%%%%%%%%%%%%%%%%%%%%%%%%%%%%%%%%%%%%%%%%%%%%%%%%%%%%%%%%%
% SECTION 3: Exercise 1 -- PING
%%%%%%%%%%%%%%%%%%%%%%%%%%%%%%%%%%%%%%%%%%%%%%%%%%%%%%%%%%%%%%%%
%%%%%%%%%%%%%%%%%%%%%%%%%%%%%%%%%%%%%%%%%%%%%%%%%%%%%%%%%%%%%%%%

\section{Exercise 1: PING}

\begin{frame}
\frametitle{Exercise 1: PING the Ensembl Server}

\textbf{File}: \texttt{P07/e1.py}

\vspace{0.3cm}

\textbf{Goal}: check whether the Ensembl service is \textbf{alive} using the \texttt{info/ping} endpoint.

\vspace{0.3cm}

\textbf{Endpoint URL}:
\begin{center}
\small\texttt{https://rest.ensembl.org/info/ping?content-type=application/json}
\end{center}

\vspace{0.3cm}

The URL has three parts:
\begin{itemize}
\item \textbf{Server}: \texttt{rest.ensembl.org}
\item \textbf{Endpoint}: \texttt{/info/ping}
\item \textbf{Parameters}: \texttt{content-type=application/json}
\end{itemize}

\end{frame}

\begin{frame}[fragile]
\frametitle{Exercise 1: Implementation}

Use the \pythoninline{http.client} module to connect to the server:

\vspace{0.2cm}

\begin{python}
import http.client
import json

SERVER = "rest.ensembl.org"
ENDPOINT = "/info/ping"
PARAMS = "?content-type=application/json"

conn = http.client.HTTPSConnection(SERVER)
conn.request("GET", ENDPOINT + PARAMS)

response = conn.getresponse()
data = json.loads(response.read().decode())

if data["ping"] == 1:
    print("ALIVE!")
\end{python}

\vspace{0.2cm}

Print the \textbf{server name}, \textbf{URL} and the \textbf{status}.

\end{frame}

%%%%%%%%%%%%%%%%%%%%%%%%%%%%%%%%%%%%%%%%%%%%%%%%%%%%%%%%%%%%%%%%
%%%%%%%%%%%%%%%%%%%%%%%%%%%%%%%%%%%%%%%%%%%%%%%%%%%%%%%%%%%%%%%%
% SECTION 4: Exercise 2 -- Gene Identification
%%%%%%%%%%%%%%%%%%%%%%%%%%%%%%%%%%%%%%%%%%%%%%%%%%%%%%%%%%%%%%%%
%%%%%%%%%%%%%%%%%%%%%%%%%%%%%%%%%%%%%%%%%%%%%%%%%%%%%%%%%%%%%%%%

\section{Exercise 2: Gene Identification}

\begin{frame}[fragile]
\frametitle{Exercise 2: Gene Identification Strings}

\textbf{File}: \texttt{P07/e2.py}

\vspace{0.3cm}

Every sequence in Ensembl has a \textbf{stable identifier} (id).

\vspace{0.1cm}
Example: the \textbf{SRCAP} gene has id \texttt{ENSG00000080603}.

\vspace{0.3cm}

\textbf{Task}: find the stable identifiers for all 10 genes and create a \textbf{dictionary}:

\vspace{0.2cm}

\begin{python}
genes = {
    "FRAT1": "ENSG00000165879",
    "ADA":   "ENSG00000196839",
    "FXN":   "ENSG00000165060",
    # ... rest of the genes
}
\end{python}

\vspace{0.2cm}

Print all entries: \pythoninline{genes["FRAT1"]} $\rightarrow$ \texttt{ENSG00000165879}

\end{frame}

%%%%%%%%%%%%%%%%%%%%%%%%%%%%%%%%%%%%%%%%%%%%%%%%%%%%%%%%%%%%%%%%
%%%%%%%%%%%%%%%%%%%%%%%%%%%%%%%%%%%%%%%%%%%%%%%%%%%%%%%%%%%%%%%%
% SECTION 5: Exercise 3 -- Getting a Remote Gene
%%%%%%%%%%%%%%%%%%%%%%%%%%%%%%%%%%%%%%%%%%%%%%%%%%%%%%%%%%%%%%%%
%%%%%%%%%%%%%%%%%%%%%%%%%%%%%%%%%%%%%%%%%%%%%%%%%%%%%%%%%%%%%%%%

\section{Exercise 3: Getting a Remote Gene}

\begin{frame}
\frametitle{Exercise 3: Getting a Remote Gene}

\textbf{File}: \texttt{P07/e3.py}

\vspace{0.3cm}

\textbf{Goal}: connect to the \texttt{/sequence/id} endpoint and get the \textbf{sequence} and \textbf{description} of the \textbf{MIR633} gene.

\vspace{0.3cm}

\textbf{Endpoint documentation}:
\begin{center}
\small\url{https://rest.ensembl.org/documentation/info/sequence_id}
\end{center}

\vspace{0.3cm}

The program should print:
\begin{itemize}
\item The gene \textbf{description}
\item The gene \textbf{sequence}
\end{itemize}

\end{frame}

\begin{frame}[fragile]
\frametitle{Exercise 3: Accessing the Endpoint}

\begin{python}
ENDPOINT = "/sequence/id/"
GENE_ID = "ENSG00000207552"  # MIR633
PARAMS = "?content-type=application/json"

conn = http.client.HTTPSConnection(SERVER)
conn.request("GET",
    ENDPOINT + GENE_ID + PARAMS)

response = conn.getresponse()
data = json.loads(response.read().decode())

print("Description:", data["desc"])
print("Sequence:", data["seq"])
\end{python}

\vspace{0.2cm}

The response JSON object contains fields like \texttt{desc}, \texttt{seq}, \texttt{id}, etc.

\end{frame}

%%%%%%%%%%%%%%%%%%%%%%%%%%%%%%%%%%%%%%%%%%%%%%%%%%%%%%%%%%%%%%%%
%%%%%%%%%%%%%%%%%%%%%%%%%%%%%%%%%%%%%%%%%%%%%%%%%%%%%%%%%%%%%%%%
% SECTION 6: Exercise 4 -- Processing Gene Information
%%%%%%%%%%%%%%%%%%%%%%%%%%%%%%%%%%%%%%%%%%%%%%%%%%%%%%%%%%%%%%%%
%%%%%%%%%%%%%%%%%%%%%%%%%%%%%%%%%%%%%%%%%%%%%%%%%%%%%%%%%%%%%%%%

\section{Exercise 4: Processing Gene Information}

\begin{frame}
\frametitle{Exercise 4: Processing Gene Information}

\textbf{File}: \texttt{P07/e4.py}

\vspace{0.3cm}

\textbf{Goal}: ask the user for a \textbf{gene name}, retrieve it from the server, and calculate:

\vspace{0.2cm}

\begin{itemize}
\item Gene \textbf{name} and \textbf{description}
\item \textbf{Total length} of the sequence
\item \textbf{Number of bases} (A, C, G, T)
\item \textbf{Percentage} of each base
\item The \textbf{most frequent base}
\end{itemize}

\vspace{0.3cm}

Use the \textbf{Seq Class} from previous assignments for the calculations.

\end{frame}

\begin{frame}[fragile]
\frametitle{Exercise 4: Expected Output}

\begin{python}[language={}]
Enter the gene name: FRAT1

Gene: FRAT1
Description: FRAT regulator of WNT ...
Total length: 1801
A: 432 (24.0%)
C: 525 (29.2%)
G: 475 (26.4%)
T: 369 (20.5%)
Most frequent base: C
\end{python}

\vspace{0.2cm}

\textbf{Steps}:
\begin{enumerate}
\item Read gene name from the user
\item Look up its \textbf{stable id} in the dictionary
\item Connect to the \texttt{/sequence/id} endpoint
\item Parse the JSON response and perform calculations
\end{enumerate}

\end{frame}

%%%%%%%%%%%%%%%%%%%%%%%%%%%%%%%%%%%%%%%%%%%%%%%%%%%%%%%%%%%%%%%%
%%%%%%%%%%%%%%%%%%%%%%%%%%%%%%%%%%%%%%%%%%%%%%%%%%%%%%%%%%%%%%%%
% SECTION 7: Exercise 5 -- All Genes
%%%%%%%%%%%%%%%%%%%%%%%%%%%%%%%%%%%%%%%%%%%%%%%%%%%%%%%%%%%%%%%%
%%%%%%%%%%%%%%%%%%%%%%%%%%%%%%%%%%%%%%%%%%%%%%%%%%%%%%%%%%%%%%%%

\section{Exercise 5: Information of All Genes}

\begin{frame}
\frametitle{Exercise 5: Information of All Genes}

\textbf{File}: \texttt{P07/e5.py}

\vspace{0.3cm}

\textbf{Goal}: perform the \textbf{same calculations} as Exercise 4 but for \textbf{all 10 genes} automatically.

\vspace{0.3cm}

\begin{itemize}
\item Gene names should be obtained from the \textbf{dictionary} (not from the user)
\item Iterate over all entries and connect to the server for each one
\item Print the results for \textbf{every gene}
\end{itemize}

\vspace{0.3cm}

\textbf{Hint}: loop through the dictionary:

\vspace{0.1cm}

\begin{center}
\pythoninline{for gene_name in genes:}
\end{center}

\end{frame}

%%%%%%%%%%%%%%%%%%%%%%%%%%%%%%%%%%%%%%%%%%%%%%%%%%%%%%%%%%%%%%%%
%%%%%%%%%%%%%%%%%%%%%%%%%%%%%%%%%%%%%%%%%%%%%%%%%%%%%%%%%%%%%%%%
% SECTION 8: Deliverables
%%%%%%%%%%%%%%%%%%%%%%%%%%%%%%%%%%%%%%%%%%%%%%%%%%%%%%%%%%%%%%%%
%%%%%%%%%%%%%%%%%%%%%%%%%%%%%%%%%%%%%%%%%%%%%%%%%%%%%%%%%%%%%%%%

\section{Deliverables}

\begin{frame}
\frametitle{Deliverables}

\textbf{P07 folder should contain}:

\vspace{0.2cm}

\begingroup\small
\begin{itemize}
\item \texttt{e1.py} -- Ping the Ensembl server
\item \texttt{e2.py} -- Gene identification dictionary
\item \texttt{e3.py} -- Get remote gene (MIR633)
\item \texttt{e4.py} -- Process a single gene (user input)
\item \texttt{e5.py} -- Process all 10 genes
\end{itemize}
\endgroup

\vspace{0.3cm}

\textbf{Remember}: Push everything to your remote Github repository!

\end{frame}

%%%%%%%%%%%%%%%%%%%%%%%%%%%%%%%%%%%%%%%%%%%%%%%%%%%%%%%%%%%%%%%%
%%%%%%%%%%%%%%%%%%%%%%%%%%%%%%%%%%%%%%%%%%%%%%%%%%%%%%%%%%%%%%%%
% SECTION 9: End
%%%%%%%%%%%%%%%%%%%%%%%%%%%%%%%%%%%%%%%%%%%%%%%%%%%%%%%%%%%%%%%%
%%%%%%%%%%%%%%%%%%%%%%%%%%%%%%%%%%%%%%%%%%%%%%%%%%%%%%%%%%%%%%%%

\begin{frame}
\begin{center}
\Huge{Questions?}

\vspace{1cm}

\Large{Time to code!}
\end{center}
\end{frame}

\end{document}
